% 第九卷 6.2.3.2 节：分页键值缓存的块表。块大小取 4 个词元（教学示意）。
% 请求 A、B 共享 8 个词元的系统提示 S1–S8（逻辑块 0、1 → 物理块 7、1，引用计数 2）；
% A 的逻辑块 2 → 物理块 3（只用 1 个槽位），B 的逻辑块 2 → 物理块 5（用 3 个槽位）。
\begin{tikzpicture}[
  x=1cm,y=1cm,
  every node/.style={font=\figurefont\scriptsize,text=FigureInk},
  slot/.style={rectangle,draw=white,line width=0.5pt,minimum width=4.6mm,minimum height=4.6mm,inner sep=0pt},
  sh/.style={slot,fill=FigureBlueFill},
  ao/.style={slot,fill=FigureYellowFill},
  bo/.style={slot,fill=FigureRedFill},
  em/.style={slot,fill=FigurePaper},
  fr/.style={rectangle,draw=FigureGrid,dashed,line width=0.5pt,minimum width=18.4mm,minimum height=4.6mm,inner sep=0pt,fill=white},
  frame/.style={draw=FigureStroke,line width=0.6pt},
  cell/.style={rectangle,draw=FigureStroke,line width=0.5pt,minimum width=6mm,minimum height=4.6mm,inner sep=0pt,fill=white},
  hd/.style={font=\figurefont\bfseries\footnotesize,text=FigureInk},
  lab/.style={font=\figurefont\scriptsize,text=FigureMuted}
]
  \def\w{0.46}
  % 逻辑视图。
  \node[hd,anchor=west] at (-0.9,2.55) {逻辑块（每个请求）};
  \foreach \r/\y/\nm in {A/1.2/请求$A$,B/-0.9/请求$B$} {
    \node[hd,anchor=east] at (-0.15,\y) {\nm};
    \foreach \b in {0,1,2} {
      \node[lab,anchor=south] at ({\b*2.1+2*\w},{\y+0.3}) {逻辑块\b};
    }
  }
  % 共享前缀 S1–S8。
  \foreach \y in {1.2,-0.9} {
    \foreach \i in {1,...,4} {\node[sh] at ({(\i-0.5)*\w},\y) {S\i};}
    \foreach \i in {5,...,8} {\node[sh] at ({2.1+(\i-4.5)*\w},\y) {S\i};}
  }
  \node[ao] at ({4.2+0.5*\w},1.2) {a1};
  \foreach \i in {2,3,4} {\node[em] at ({4.2+(\i-0.5)*\w},1.2) {};}
  \foreach \i in {1,2,3} {\node[bo] at ({4.2+(\i-0.5)*\w},-0.9) {b\i};}
  \node[em] at ({4.2+3.5*\w},-0.9) {};
  \foreach \y in {1.2,-0.9} {\foreach \b in {0,1,2} {\draw[frame] ({\b*2.1},{\y-0.23}) rectangle ++({4*\w},0.46);}}

  % 块表。
  \node[hd] at (7.55,2.55) {块表};
  \foreach \y in {1.2,-0.9} {\node[lab,anchor=south] at (7.55,{\y+0.72}) {逻辑$\to$物理};}
  \foreach \k/\v in {0/7,1/1,2/3} {\node[cell] (TA\k) at (7.55,{1.2+0.46-\k*0.46}) {\k$\to$\v};}
  \foreach \k/\v in {0/7,1/1,2/5} {\node[cell] (TB\k) at (7.55,{-0.9+0.46-\k*0.46}) {\k$\to$\v};}
  \draw[-{Stealth[length=1.6mm]},FigureMuted,line width=0.5pt] (6.15,1.2) -- (TA1.west);
  \draw[-{Stealth[length=1.6mm]},FigureMuted,line width=0.5pt] (6.15,-0.9) -- (TB1.west);

  % 物理块池（自上而下为物理块 0–7）。
  \node[hd,anchor=west] at (9.75,2.55) {物理块（显存）};
  \foreach \p in {0,2,4,6} {\node[fr,anchor=west] at (9.8,{1.75-\p*0.56}) {};}
  \foreach \p in {0,2,4,6} {\node[lab,anchor=west] at (11.75,{1.75-\p*0.56}) {块\p：空闲};}
  % 物理块 1：S5–S8；物理块 7：S1–S4。
  \foreach \i in {5,...,8} {\node[sh] at ({9.8+(\i-4.5)*\w},{1.75-0.56}) {S\i};}
  \foreach \i in {1,...,4} {\node[sh] at ({9.8+(\i-0.5)*\w},{1.75-7*0.56}) {S\i};}
  \node[ao] at ({9.8+0.5*\w},{1.75-3*0.56}) {a1};
  \foreach \i in {2,3,4} {\node[em] at ({9.8+(\i-0.5)*\w},{1.75-3*0.56}) {};}
  \foreach \i in {1,2,3} {\node[bo] at ({9.8+(\i-0.5)*\w},{1.75-5*0.56}) {b\i};}
  \node[em] at ({9.8+3.5*\w},{1.75-5*0.56}) {};
  \foreach \p in {1,3,5,7} {\draw[frame] (9.8,{1.75-\p*0.56-0.23}) rectangle ++({4*\w},0.46);}
  \node[lab,anchor=west] at (11.75,{1.75-0.56}) {块1：引用计数2};
  \node[lab,anchor=west] at (11.75,{1.75-7*0.56}) {块7：引用计数2};
  \node[lab,anchor=west] at (11.75,{1.75-3*0.56}) {块3：引用计数1};
  \node[lab,anchor=west] at (11.75,{1.75-5*0.56}) {块5：引用计数1};

  % 块表到物理块的映射。
  \foreach \t/\p in {TA0/7,TA1/1,TA2/3,TB0/7,TB1/1,TB2/5} {
    \draw[-{Stealth[length=1.6mm]},FigureStroke,line width=0.5pt] (\t.east) -- (9.78,{1.75-\p*0.56});
  }
\end{tikzpicture}
